\section{复杂度和暴力基线章正文案例推导补强}

这一节补齐本章正文出现过的 LeetCode 案例推导。阅读顺序统一按“题目说明、样例入口、暴力基线、暴力过程、重复工作、优化条件、相邻状态、状态复用、指针和字段、代码落点、正确性、边界实验、复杂度变化、迁移追问”展开；目的不是新增题量，而是把正文中的案例都讲成可复述、可证明、可迁移的解题链。

\subsection{LeetCode 1 两数之和}

\textbf{题目说明。} LeetCode 1 两数之和 的题面入口要先还原成具体任务：当前元素只需要寻找唯一补数，历史值到下标的映射就是最小状态。

\textbf{样例入口。} 拿数组 [2, 7]、target=9 手算：站在 2 时历史为空，不能配对；站在 7 时只需要问历史里有没有 2。再试 [3, 3]、target=6：如果先把当前 3 插入再查，就会把同一个下标用两次。

\textbf{暴力基线。} 暴力两层循环枚举 i 和 j，检查 nums[i]+nums[j] 是否等于 target，找到后返回两个下标。

\textbf{暴力过程。} 以 nums=[2, 7, 11, 15]、target=9 为例，暴力先试 2+7 命中；换成 [3, 2, 4]、target=6 时，会试 3+2 失败，再试 2+4 命中。

\textbf{重复工作。} 题面模型：当前元素只需要寻找唯一补数，历史值到下标的映射就是最小状态。暴力版的慢点要落到本题：浪费在每个当前位置都回头扫描历史。站在 nums[i] 时，真正需要的不是所有历史值，而是历史里有没有 target-nums[i]。后面的优化只消掉这类重复工作，不能改变答案集合。

\textbf{优化条件。} 本题的关键观察是：先查再插入，避免同一个下标被用两次；若数组有序才考虑双指针。这句话必须能翻译成可维护状态；如果题目条件改变后观察不再成立，就不能继续套原来的写法。

\textbf{相邻状态。} 规律不是“用哈希表”四个字，而是每个当前位置只和它左侧历史配对，代码必须先查补数再插入当前值。把这个特例继续拆开看：每个合法答案都要还原成当前状态和某个历史状态的配对；哈希表的 key 负责定位历史状态，value 负责表达次数、最早位置或存在性。先查还是先写，必须由是否允许当前前缀和自己配对来决定。

\textbf{状态复用。} 哈希表保存已经扫描过的值到下标的映射。每到一个新数，先查补数是否存在；不存在再把当前数放入历史。手算时只改被新输入影响的那一小部分，而不是回头重算完整候选。

\textbf{指针和字段。} 状态字段要能说清：i 是当前下标，need=target-nums[i] 是补数，seen[value] 是该值在左侧历史中出现的下标。手算时按这个协议跟踪：nums=[3, 3]、target=6 时，第一个 3 先查不到补数，再插入；第二个 3 查到历史 3，返回两个不同下标。

\textbf{代码落点。} 关键顺序是先 find(need)，再 seen[nums[i]]=i。若先插入当前值，target=2*nums[i] 时可能把同一位置用两次。关键代码行必须能逐句翻译成“旧状态怎样变成新状态”，否则就只是模板。

\textbf{正确性。} 每个答案对都有一个较右下标；扫描到较右下标时，较左下标已经在 seen 中，所以一定能被查到。

\textbf{边界实验。} 先用重复数字、目标为偶数且补数等于当前值、没有答案时返回空结果做反例检查。实验时把重复数字、目标为偶数且补数等于当前值、没有答案时返回空结果加入对拍集合，用平方暴力和优化版比较。失败时先看初始历史状态、查询更新顺序和哈希表 value 类型。

\textbf{复杂度变化。} 暴力 O(\ensuremath{n^{2}})，哈希表扫描 O(n) 均摊时间、O(n) 空间。

\textbf{迁移追问。} 如果题目改成“若改成返回所有不重复数对，要先排序或用频次表处理重复”，先判断原状态是否仍然覆盖新答案集合，再决定是微调边界、改 value 语义，还是换算法。

\subsection{LeetCode 560 和为 K 的子数组}

\textbf{题目说明。} LeetCode 560 和为 K 的子数组给定整数数组 nums 和整数 k，统计和等于 k 的连续子数组个数；数组中可以有负数。

\textbf{样例入口。} 样例 nums=[1, 1, 1]，k=2，输出 2；两个合法连续子数组分别是前两个 1 和后两个 1。

\textbf{暴力基线。} 暴力固定 left，向右累计 sum；每当 sum==k 就计数。

\textbf{暴力过程。} 以 nums=[1, 1, 1]、k=2 为例，从 left=0 可找到 [1, 1]，从 left=1 又找到后两个 1。

\textbf{重复工作。} 题面模型：当前前缀和减去目标值，就是需要匹配的历史前缀。暴力版的慢点要落到本题：浪费在每个右端都回头枚举所有 left。区间和等于 k 可以改写成 prefix[right]-prefix[left]=k。后面的优化只消掉这类重复工作，不能改变答案集合。

\textbf{优化条件。} 本题的关键观察是：哈希表保存前缀和出现次数，先查询再更新。这句话必须能翻译成可维护状态；如果题目条件改变后观察不再成立，就不能继续套原来的写法。

\textbf{相邻状态。} 规律是哈希表保存前缀和频次，且初始要放 0:1。把这个特例继续拆开看：每个合法答案都要还原成当前状态和某个历史状态的配对；哈希表的 key 负责定位历史状态，value 负责表达次数、最早位置或存在性。先查还是先写，必须由是否允许当前前缀和自己配对来决定。

\textbf{状态复用。} 扫描前缀和 prefix。哈希表 count[old\_prefix] 保存历史前缀出现次数；当前位置需要查 prefix-k 出现过几次。手算时只改被新输入影响的那一小部分，而不是回头重算完整候选。

\textbf{指针和字段。} 状态字段要能说清：prefix 是当前前缀和，need=prefix-k 是所需历史前缀，count 保存历史前缀频次，answer 累加匹配次数。手算时按这个协议跟踪：[1, 1, 1] 中扫描到第二个元素时 prefix=2，需要历史 0，count[0]=1，所以找到第一个长度 2 子数组。

\textbf{代码落点。} 关键是先 answer += count[prefix-k]，再 count[prefix]++。初始化 count[0]=1，表示从下标 0 开始的区间。关键代码行必须能逐句翻译成“旧状态怎样变成新状态”，否则就只是模板。

\textbf{正确性。} 任意和为 k 的子数组都对应一对前缀差 k；扫描到右端时，所有合法左端前缀都已在表里。

\textbf{边界实验。} 先用负数、目标为零、从下标零开始的区间、重复前缀做反例检查。实验时把负数、目标为零、从下标零开始的区间、重复前缀加入对拍集合，用平方暴力和优化版比较。失败时先看初始历史状态、查询更新顺序和哈希表 value 类型。

\textbf{复杂度变化。} 暴力 O(\ensuremath{n^{2}})，前缀哈希 O(n) 均摊时间、O(n) 空间。

\textbf{迁移追问。} 如果题目改成“若要求最长区间，哈希表应保存最早位置而不是次数”，先判断原状态是否仍然覆盖新答案集合，再决定是微调边界、改 value 语义，还是换算法。

\subsection{LeetCode 70 爬楼梯}

\textbf{题目说明。} LeetCode 70 爬楼梯 的题面入口要先还原成具体任务：最后一步来自前一阶或前两阶，是最小的重叠子问题模型。

\textbf{样例入口。} 拿 n=4 手算，最后一步要么从第 3 阶跨 1 阶，要么从第 2 阶跨 2 阶，这两类互斥且覆盖所有走法。于是 ways[4]=ways[3]+ways[2]。

\textbf{暴力基线。} 暴力递归枚举每一步走 1 阶或 2 阶，直到正好到 n 计数，超过 n 则失败。

\textbf{暴力过程。} 以 n=4 为例，最后一步要么从第 3 阶走 1 阶，要么从第 2 阶走 2 阶；这两类互斥且覆盖所有走法。

\textbf{重复工作。} 题面模型：最后一步来自前一阶或前两阶，是最小的重叠子问题模型。暴力版的慢点要落到本题：浪费在递归树反复计算 ways(2)、ways(3) 等同一子问题。到达某一阶后，后面的走法数量和之前怎么来无关。后面的优化只消掉这类重复工作，不能改变答案集合。

\textbf{优化条件。} 本题的关键观察是：滚动变量保存最近两个状态即可，不需要完整数组。这句话必须能翻译成可维护状态；如果题目条件改变后观察不再成立，就不能继续套原来的写法。

\textbf{相邻状态。} 规律不是背斐波那契，而是最后一步只有两个来源，状态只需保存前两项。把这个特例继续拆开看：先问暴力搜索里哪些不同路径到达同一个后续问题，再给这个后续问题命名为状态。状态必须包含未来决策需要的全部信息，转移只从更小且已经算清的状态来。

\textbf{状态复用。} dp[i] 表示到达第 i 阶的方法数，转移是 dp[i]=dp[i-1]+dp[i-2]。再压缩成两个滚动变量。手算时只改被新输入影响的那一小部分，而不是回头重算完整候选。

\textbf{指针和字段。} 状态字段要能说清：prev2 表示 dp[i-2]，prev1 表示 dp[i-1]，current 表示 dp[i]。手算时按这个协议跟踪：n=4 时 dp[1]=1、dp[2]=2、dp[3]=3、dp[4]=5。每一步只依赖前两项。

\textbf{代码落点。} 关键是初始化：n=0 按题意通常返回 1 或不出现；LeetCode 输入 n>=1，n=1 返回 1，n=2 返回 2。关键代码行必须能逐句翻译成“旧状态怎样变成新状态”，否则就只是模板。

\textbf{正确性。} 所有到第 i 阶的方案按最后一步分成从 i-1 来和从 i-2 来两类，二者不重不漏。

\textbf{边界实验。} 先用n 为一、n 为二、题目是否允许零阶、结果范围做反例检查。实验时覆盖n 为一、n 为二、题目是否允许零阶、结果范围，先打印完整 DP 表，再比较空间压缩版本。若压缩版不同，优先检查遍历顺序和旧值保存。

\textbf{复杂度变化。} 暴力递归指数级；DP O(n) 时间，滚动变量 O(1) 空间。

\textbf{迁移追问。} 如果题目改成“若每次可走一到 k 阶，转移变成固定宽度窗口求和”，先判断原状态是否仍然覆盖新答案集合，再决定是微调边界、改 value 语义，还是换算法。

\subsection{LeetCode 1011 在 D 天内送达包裹}

\textbf{题目说明。} LeetCode 1011 在 D 天内送达包裹 的题面入口要先还原成具体任务：容量越大，需要天数越少，存在最小可行容量。

\textbf{样例入口。} 拿 weights=[1, 2, 3, 4, 5]、D=3 手算，容量 6 可以分成 [1, 2, 3]、[4]、[5] 三天；容量 5 需要更多天。

\textbf{暴力基线。} 暴力从运载能力 1 开始试到 sum(weights)，每个能力模拟装船天数，找到第一个能在 days 内完成的能力。

\textbf{暴力过程。} weights=[1, 2, 3, 4, 5]、days=3 时，capacity=6 可分成 [1, 2, 3]、[4]、[5] 三天；capacity=5 则需要更多天。

\textbf{重复工作。} 题面模型：容量越大，需要天数越少，存在最小可行容量。暴力版的慢点要落到本题：浪费在能力值有单调性却线性试。运载能力越大，需要天数只会不增。后面的优化只消掉这类重复工作，不能改变答案集合。

\textbf{优化条件。} 本题的关键观察是：下界是最大包裹，上界是总重量，检查函数按顺序贪心装船。这句话必须能翻译成可维护状态；如果题目条件改变后观察不再成立，就不能继续套原来的写法。

\textbf{相邻状态。} 规律是容量越大所需天数越少，可行性单调，二分最小容量。把这个特例继续拆开看：每次比较 mid 之后，必须能指出哪一侧整体不可能包含目标，或哪一侧整体已经满足可行性。二分的核心不是 mid 公式，而是被丢弃区间确实不含答案。

\textbf{状态复用。} 答案范围是 [max(weights), sum(weights)]。二分 capacity，check(capacity) 贪心从左到右装，超载就开新一天。手算时只改被新输入影响的那一小部分，而不是回头重算完整候选。

\textbf{指针和字段。} 状态字段要能说清：capacity 是候选船载重，used\_days 是模拟所需天数，current\_load 是当天已装重量。手算时按这个协议跟踪：capacity=6 时，1+2+3 正好一船；再装 4 会超载，于是开第二天；最后 5 第三天，满足 days=3。

\textbf{代码落点。} 关键是 left 从最大单件重量开始，不能从 1 开始；check 返回 used\_days<=days 时，说明能力可行，right=mid。关键代码行必须能逐句翻译成“旧状态怎样变成新状态”，否则就只是模板。

\textbf{正确性。} 固定 capacity 时，尽量把当天装满的贪心会让天数最少；可行性随 capacity 增大单调变真，二分第一个 true。

\textbf{边界实验。} 先用D 为一、D 等于包裹数、单个包裹超过猜测容量做反例检查。实验时覆盖D 为一、D 等于包裹数、单个包裹超过猜测容量，逐轮打印 left、mid、right、比较值或检查结果。只要有一次被排除区间仍含答案，二分不变量就错了。

\textbf{复杂度变化。} 线性试能力 O(sum*n)，二分 O(n log(sum-max)) 时间，O(1) 空间。

\textbf{迁移追问。} 如果题目改成“它和分割数组最大值都是连续分段最大和模型”，先判断原状态是否仍然覆盖新答案集合，再决定是微调边界、改 value 语义，还是换算法。

\begin{principlebox}{本节复盘方式}
不要只看最终算法名。每道题复盘时先遮住优化版，口述暴力枚举对象；再指出相邻窗口、历史前缀、候选区间、搜索状态或子问题里哪一部分被重复处理；最后用一个边界样例验证状态更新顺序。能讲完这三步，才算真正掌握本章案例。
\end{principlebox}
